\subsection{Contribution-isolating comparisons}
\label{sec:r15baselines}

The comparison separates continuation regression, direct policy optimization,
and direct approximation of the Bellman equation.  The direct HJB method uses
the same two learned scalar components as NBO: a time network of widths
$(16,16)$ and a state--time network of widths $(32,32)$, each with hyperbolic
tangent activations and a scalar output.  Its known lift is the full-value
lift, so the terminal condition is imposed exactly.  NBO's known lift instead
represents its postdecision continuation target.  This difference follows from
the objects being approximated; the learned residual capacity is the same.
Both methods use the prespecified occupation sample and the same available
model information.  No reference solution enters either training objective.

Write $c=a-(\sigma_I^2+\sigma_C^2)/2$ and
$\Sigma\Sigma'=\sigma_I^2I+\sigma_C^2\mathbf1\mathbf1'$.  The direct method
minimizes the squared residual of
\begin{equation}
 \begin{aligned}
 0={}&-v_t-\max_{m\in\mathcal M_t}
 \{u(y,m)+(c\mathbf1+\kappa\tanh(By)-m)'D_yv\}\\
 &-\tfrac12\operatorname{tr}(\Sigma\Sigma'D_y^2v)+\rho v,
 \qquad v(T,y)=g(y),
 \end{aligned}
 \label{eq:r15-direct-hjb}
\end{equation}
where $\mathcal M_t$ is the action tube shared by the four primary methods.
For $q=dD_yv$, the action-dependent Hamiltonian is
\begin{equation}
 H_q(m)=\frac1d\sum_i\log m_i-\frac\zeta2\bar m^2
          -\frac1d\sum_iq_i m_i.
 \label{eq:r15-hjb-hamiltonian}
\end{equation}
It is strictly concave.  Its maximizer is obtained from the scalar consistency
equation $z=d^{-1}\sum_i m_i(z)$, with
$m_i(z)=\operatorname{proj}_{[\ell_t,u_t]}(q_i+\zeta z)^{-1}$ when
$q_i+\zeta z>0$, and $m_i(z)=u_t$ otherwise.  The implementation uses the same
44 monotone bisection steps as the other feasible costate comparisons.
The direct HJB policy evaluates a critic gradient and this maximization at
every decision.  NBO deploys its separately learned actor.  The resulting
difference in online computation is included in validation, stopping,
confirmation, and deployment costs.

Two independent Gaussian probe banks estimate the covariance Hessian trace.
If $\widehat R_1$ and $\widehat R_2$ are the corresponding residual estimates,
then, conditional on a collocation state and the current critic,
\begin{equation}
 \mathbb E[\widehat R_1\widehat R_2]=R[v]^2.
 \label{eq:r15-hjb-independent-traces}
\end{equation}
The product, rather than the square of one estimate, is the training loss.
Its realized value can be negative.  Negative batches are retained.  This
identity specifies an optimization objective; the common payoff verifier
supplies the economic comparison.

\paragraph*{Classical scalar training.}
The original volatility calibration also has a scalar comparison.  A monotone
implicit Howard method solves the one-state problem and saves an action at
every time and state node.  These saved policies are deployed by interpolation
on fresh diffusion paths, together with the analytical schedule and the
hash-pinned scalar NBO and DPO policies from the earlier experiment.  Full-box
and matched-tube classical policies are separate candidates.  Three successively
refined state--time grids are followed by an enlargement of the state interval
at fixed spatial and temporal spacing.  The deployment meshes aggregate the
same fine Brownian increments.  Policy differences are reported both over
the declared initial population and at its individual states $-1$, $0$, and
$1$.  Historical neural fitting clocks retain their historical interpretation;
the new classical solves and all new deployment work are measured separately.

The scalar reference has a numerical error account at the level of the
declared finite-grid equations.  Its algebraic error account uses an outward
residual calculation.  State-domain approximation and diffusion time
discretization remain separate from this algebraic account.  The scalar
comparison supplies an independent low-dimensional accuracy and work
benchmark; the high-dimensional method comparison uses its own payoff data.
